\documentclass[11pt]{article}
% ---------------------------------------------------------------
% R4.02 Qualite de developpement -- preambule commun des sujets de TP et des corrections
% Article clair, fond blanc, accents de couleur coherents avec les slides
% Charte reprise de R3.07 (meme palette IUT que les diapositives)
% Utilisation :
%   \documentclass[11pt]{article}
%   % ---------------------------------------------------------------
% R4.02 Qualite de developpement -- preambule commun des sujets de TP et des corrections
% Article clair, fond blanc, accents de couleur coherents avec les slides
% Charte reprise de R3.07 (meme palette IUT que les diapositives)
% Utilisation :
%   \documentclass[11pt]{article}
%   % ---------------------------------------------------------------
% R4.02 Qualite de developpement -- preambule commun des sujets de TP et des corrections
% Article clair, fond blanc, accents de couleur coherents avec les slides
% Charte reprise de R3.07 (meme palette IUT que les diapositives)
% Utilisation :
%   \documentclass[11pt]{article}
%   \input{../../latex/common/r402-tp-preamble.tex}
%   \begin{document}
%   \rtpheader{1}{Titre du TP}
% ---------------------------------------------------------------
\usepackage[T1]{fontenc}
\usepackage[utf8]{inputenc}
\usepackage[french]{babel}
\usepackage[scale=0.92]{tgheros}
\usepackage{tgpagella}
\usepackage[a4paper, margin=2.2cm, headheight=14pt]{geometry}
\usepackage{graphicx}
\usepackage{xcolor}
\usepackage{listings}
\usepackage[skins,listings,breakable]{tcolorbox}
\usepackage{enumitem}
\usepackage{booktabs}
\usepackage{array}
\usepackage{multirow}
\usepackage{tikz}
\usetikzlibrary{shapes, shapes.geometric, arrows, arrows.meta, positioning, calc}
\usepackage[colorlinks=true, linkcolor=qaaccentdark, urlcolor=qaaccentdark]{hyperref}
\usepackage{fancyhdr}
\usepackage{lastpage}
\usepackage{titlesec}
\usepackage{etoolbox}
\usepackage{eurosym}

% ---------------------------------------------------------------
% MENTION D'ELABORATION (discrete : gris clair, petit corps)
% ---------------------------------------------------------------
\definecolor{qaaimention}{HTML}{A3AAB5}
\newcommand{\raimention}{{\scriptsize\color{qaaimention}Supports \'elabor\'es avec l'aide de Claude Code Opus 5.5}}

% ---------------------------------------------------------------
% PALETTE -- version imprimable des couleurs des diapositives
% ---------------------------------------------------------------
\definecolor{qaaccentdark}{HTML}{7A1680}   % violet du logo IUT, fonce pour rester lisible sur blanc
\definecolor{qaaccent2dark}{HTML}{9B4B94}  % violet clair du logo IUT (accent secondaire)
\definecolor{qapanellight}{HTML}{F4F6F8}   % fond de bloc tres clair
\definecolor{qalinelight}{HTML}{D8DEE4}    % filets
\definecolor{qatextdark}{HTML}{1A1D23}     % texte principal
\definecolor{qamuteddark}{HTML}{5B6472}    % texte secondaire
\definecolor{qacodekwl}{HTML}{7A1680}      % code : mots-cles
\definecolor{qacodestrl}{HTML}{2E7D32}     % code : chaines
\definecolor{qacodecoml}{HTML}{6B7280}     % code : commentaires
\definecolor{qacodeannl}{HTML}{B26A00}     % code : annotations / javadoc

\color{qatextdark}

\setlength{\parskip}{4pt}
\setlength{\parindent}{0pt}

% ---------------------------------------------------------------
% TITRES DE SECTION -- texte blanc sur bande de couleur pleine largeur
% (gere aussi les variantes etoilees \section*)
% ---------------------------------------------------------------
\newcommand{\qasectionband}[2]{%
  \noindent\colorbox{#1}{%
    \parbox[c]{\dimexpr\linewidth-2\fboxsep\relax}{%
      \color{white}#2%
    }%
  }%
}
\newcommand{\secformat}[1]{\qasectionband{qaaccentdark}{\bfseries\Large\thesection\quad #1}}
\newcommand{\secformatnn}[1]{\qasectionband{qaaccentdark}{\bfseries\Large #1}}
\newcommand{\subsecformat}[1]{\qasectionband{qaaccent2dark}{\bfseries\large\thesubsection\quad #1}}
\newcommand{\subsecformatnn}[1]{\qasectionband{qaaccent2dark}{\bfseries\large #1}}

\titleformat{\section}{\normalfont}{}{0pt}{\secformat}
\titleformat{name=\section,numberless}{\normalfont}{}{0pt}{\secformatnn}
\titleformat{\subsection}{\normalfont}{}{0pt}{\subsecformat}
\titleformat{name=\subsection,numberless}{\normalfont}{}{0pt}{\subsecformatnn}
\titleformat{\subsubsection}{\normalfont\bfseries\color{qaaccentdark}}{\thesubsubsection}{0.6em}{}
\titlespacing*{\section}{0pt}{14pt}{8pt}
\titlespacing*{\subsection}{0pt}{10pt}{6pt}
\titlespacing*{\subsubsection}{0pt}{8pt}{4pt}

% ---------------------------------------------------------------
% LISTINGS -- theme clair (lstlisting reste utilisable tel quel)
% \lstcode : police + couleur de base, a reutiliser dans les surcharges
% ---------------------------------------------------------------
\newcommand{\lstcode}{\ttfamily\color{qatextdark}}

\lstdefinelanguage{TypeScript}{
  keywords={typeof, new, true, false, catch, function, return, null, switch, var,
    if, in, while, do, else, case, break, const, let, interface, class, extends,
    implements, enum, export, import, this, super, constructor, public, private,
    protected, static, async, await, throw, throws, try, finally, void},
  ndkeywords={boolean, any, number, string, int, long, double},
  sensitive=true,
  comment=[l]{//},
  morecomment=[s]{/*}{*/},
  morestring=[b]',
  morestring=[b]"
}

\lstdefinestyle{r402code}{
  basicstyle=\lstcode\footnotesize,
  keywordstyle=\color{qacodekwl}\bfseries,
  ndkeywordstyle=\color{qaaccent2dark}\bfseries,
  stringstyle=\color{qacodestrl},
  commentstyle=\color{qacodecoml}\itshape,
  morecomment=[s][\color{qacodeannl}]{/**}{*/},
  backgroundcolor=\color{qapanellight},
  rulecolor=\color{qalinelight},
  frame=single,
  framesep=4pt,
  xleftmargin=6pt,
  xrightmargin=6pt,
  numbers=left,
  numberstyle=\tiny\color{qamuteddark},
  numbersep=8pt,
  keepspaces=true,
  columns=fullflexible,
  showstringspaces=false,
  showspaces=false,
  showtabs=false,
  breaklines=true,
  breakatwhitespace=false,
  captionpos=b,
  tabsize=2,
  escapeinside={(*@}{@*)},
  % accents UTF-8 dans le code (listings ne gere pas l'UTF-8 multi-octets)
  literate={é}{{\'e}}1 {è}{{\`e}}1 {ê}{{\^e}}1 {à}{{\`a}}1 {â}{{\^a}}1
           {ù}{{\`u}}1 {û}{{\^u}}1 {ô}{{\^o}}1 {î}{{\^i}}1 {ç}{{\c c}}1
           {É}{{\'E}}1 {À}{{\`A}}1 {€}{{\euro{}}}1 {…}{{\ldots}}1,
}
\lstset{style=r402code}

% ---------------------------------------------------------------
% BOITES : exercice, code, information
% ---------------------------------------------------------------
\tcbset{
  exobox/.style={colback=qapanellight, colframe=qaaccentdark, boxrule=0.6pt,
    arc=2pt, left=8pt, right=8pt, top=6pt, bottom=6pt,
    coltitle=white, colbacktitle=qaaccentdark, fonttitle=\bfseries,
    boxsep=2pt, breakable},
  infobox/.style={colback=white, colframe=qaaccent2dark, boxrule=0.6pt,
    leftrule=3pt, arc=0pt, left=8pt, right=8pt, top=5pt, bottom=5pt,
    coltitle=qaaccent2dark, colbacktitle=white, fonttitle=\bfseries,
    breakable},
  codebox/.style={colback=qatextdark, colframe=qalinelight, boxrule=0.4pt,
    arc=2pt, left=8pt, right=8pt, top=5pt, bottom=5pt, listing only,
    listing options={basicstyle=\ttfamily\small\color{white},
      breaklines=true, showstringspaces=false, columns=fullflexible,
      keepspaces=true, frame=none, numbers=none, backgroundcolor={}}, breakable}
}
\newtcolorbox[auto counter]{exobox}[2][]{exobox,
  title={Exercice~\thetcbcounter\ifstrempty{#2}{}{ --- #2}}, #1}
\newtcolorbox{infobox}[1][\`A retenir]{infobox, title={#1}}
\newtcblisting{codebox}{codebox}

% ---------------------------------------------------------------
% EN-TETE DE PAGE (fil d'Ariane + pagination)
% ---------------------------------------------------------------
\newcommand{\qaheadtitlevalue}{}
\newcommand{\qaheadtitle}[1]{\renewcommand{\qaheadtitlevalue}{#1}}

\pagestyle{fancy}
\fancyhf{}
\renewcommand{\headrulewidth}{0.4pt}
\fancyhead[L]{\small\color{qamuteddark}R4.02 --- \qaheadtitlevalue}
\fancyhead[R]{\small\color{qamuteddark}\thepage\,/\,\pageref{LastPage}}
\fancyfoot[C]{\footnotesize\color{qamuteddark}IUT de Vannes --- Universit\'e Bretagne Sud}
\fancypagestyle{plain}{\fancyhf{}\renewcommand{\headrulewidth}{0pt}%
  \fancyfoot[C]{\footnotesize\color{qamuteddark}IUT de Vannes --- Universit\'e Bretagne Sud}}

% ---------------------------------------------------------------
% BANDEAU DE TITRE (pas une page dediee, pour rester compact)
% \rtpheader{numero ou intitule court}{titre du TP}
% ---------------------------------------------------------------
\newcommand{\ryear}{2026/2027}
\newcommand{\rtpheader}[2]{%
  \thispagestyle{plain}%
  \begin{minipage}[t]{0.68\textwidth}
    {\small\color{qamuteddark}
      Universit\'e Bretagne Sud\\
      IUT de Vannes\\
      D\'epartement Informatique --- \ryear}
  \end{minipage}%
  \hfill
  \begin{minipage}[t]{0.28\textwidth}
    \raggedleft
    \includegraphics[width=2.6cm]{iut_logo.jpg}
  \end{minipage}

  \vspace{6mm}
  \begin{center}
    {\large\bfseries\color{qatextdark} BUT 2 --- R4.02. Qualit\'e de d\'eveloppement}\\[2mm]
    {\normalsize\color{qamuteddark} TP #1}\\[1mm]
    {\Large\bfseries\color{qaaccentdark} #2}\\[2mm]
    \raimention
  \end{center}
  \vspace{4mm}
  \qaheadtitle{#2}
}

%   \begin{document}
%   \rtpheader{1}{Titre du TP}
% ---------------------------------------------------------------
\usepackage[T1]{fontenc}
\usepackage[utf8]{inputenc}
\usepackage[french]{babel}
\usepackage[scale=0.92]{tgheros}
\usepackage{tgpagella}
\usepackage[a4paper, margin=2.2cm, headheight=14pt]{geometry}
\usepackage{graphicx}
\usepackage{xcolor}
\usepackage{listings}
\usepackage[skins,listings,breakable]{tcolorbox}
\usepackage{enumitem}
\usepackage{booktabs}
\usepackage{array}
\usepackage{multirow}
\usepackage{tikz}
\usetikzlibrary{shapes, shapes.geometric, arrows, arrows.meta, positioning, calc}
\usepackage[colorlinks=true, linkcolor=qaaccentdark, urlcolor=qaaccentdark]{hyperref}
\usepackage{fancyhdr}
\usepackage{lastpage}
\usepackage{titlesec}
\usepackage{etoolbox}
\usepackage{eurosym}

% ---------------------------------------------------------------
% MENTION D'ELABORATION (discrete : gris clair, petit corps)
% ---------------------------------------------------------------
\definecolor{qaaimention}{HTML}{A3AAB5}
\newcommand{\raimention}{{\scriptsize\color{qaaimention}Supports \'elabor\'es avec l'aide de Claude Code Opus 5.5}}

% ---------------------------------------------------------------
% PALETTE -- version imprimable des couleurs des diapositives
% ---------------------------------------------------------------
\definecolor{qaaccentdark}{HTML}{7A1680}   % violet du logo IUT, fonce pour rester lisible sur blanc
\definecolor{qaaccent2dark}{HTML}{9B4B94}  % violet clair du logo IUT (accent secondaire)
\definecolor{qapanellight}{HTML}{F4F6F8}   % fond de bloc tres clair
\definecolor{qalinelight}{HTML}{D8DEE4}    % filets
\definecolor{qatextdark}{HTML}{1A1D23}     % texte principal
\definecolor{qamuteddark}{HTML}{5B6472}    % texte secondaire
\definecolor{qacodekwl}{HTML}{7A1680}      % code : mots-cles
\definecolor{qacodestrl}{HTML}{2E7D32}     % code : chaines
\definecolor{qacodecoml}{HTML}{6B7280}     % code : commentaires
\definecolor{qacodeannl}{HTML}{B26A00}     % code : annotations / javadoc

\color{qatextdark}

\setlength{\parskip}{4pt}
\setlength{\parindent}{0pt}

% ---------------------------------------------------------------
% TITRES DE SECTION -- texte blanc sur bande de couleur pleine largeur
% (gere aussi les variantes etoilees \section*)
% ---------------------------------------------------------------
\newcommand{\qasectionband}[2]{%
  \noindent\colorbox{#1}{%
    \parbox[c]{\dimexpr\linewidth-2\fboxsep\relax}{%
      \color{white}#2%
    }%
  }%
}
\newcommand{\secformat}[1]{\qasectionband{qaaccentdark}{\bfseries\Large\thesection\quad #1}}
\newcommand{\secformatnn}[1]{\qasectionband{qaaccentdark}{\bfseries\Large #1}}
\newcommand{\subsecformat}[1]{\qasectionband{qaaccent2dark}{\bfseries\large\thesubsection\quad #1}}
\newcommand{\subsecformatnn}[1]{\qasectionband{qaaccent2dark}{\bfseries\large #1}}

\titleformat{\section}{\normalfont}{}{0pt}{\secformat}
\titleformat{name=\section,numberless}{\normalfont}{}{0pt}{\secformatnn}
\titleformat{\subsection}{\normalfont}{}{0pt}{\subsecformat}
\titleformat{name=\subsection,numberless}{\normalfont}{}{0pt}{\subsecformatnn}
\titleformat{\subsubsection}{\normalfont\bfseries\color{qaaccentdark}}{\thesubsubsection}{0.6em}{}
\titlespacing*{\section}{0pt}{14pt}{8pt}
\titlespacing*{\subsection}{0pt}{10pt}{6pt}
\titlespacing*{\subsubsection}{0pt}{8pt}{4pt}

% ---------------------------------------------------------------
% LISTINGS -- theme clair (lstlisting reste utilisable tel quel)
% \lstcode : police + couleur de base, a reutiliser dans les surcharges
% ---------------------------------------------------------------
\newcommand{\lstcode}{\ttfamily\color{qatextdark}}

\lstdefinelanguage{TypeScript}{
  keywords={typeof, new, true, false, catch, function, return, null, switch, var,
    if, in, while, do, else, case, break, const, let, interface, class, extends,
    implements, enum, export, import, this, super, constructor, public, private,
    protected, static, async, await, throw, throws, try, finally, void},
  ndkeywords={boolean, any, number, string, int, long, double},
  sensitive=true,
  comment=[l]{//},
  morecomment=[s]{/*}{*/},
  morestring=[b]',
  morestring=[b]"
}

\lstdefinestyle{r402code}{
  basicstyle=\lstcode\footnotesize,
  keywordstyle=\color{qacodekwl}\bfseries,
  ndkeywordstyle=\color{qaaccent2dark}\bfseries,
  stringstyle=\color{qacodestrl},
  commentstyle=\color{qacodecoml}\itshape,
  morecomment=[s][\color{qacodeannl}]{/**}{*/},
  backgroundcolor=\color{qapanellight},
  rulecolor=\color{qalinelight},
  frame=single,
  framesep=4pt,
  xleftmargin=6pt,
  xrightmargin=6pt,
  numbers=left,
  numberstyle=\tiny\color{qamuteddark},
  numbersep=8pt,
  keepspaces=true,
  columns=fullflexible,
  showstringspaces=false,
  showspaces=false,
  showtabs=false,
  breaklines=true,
  breakatwhitespace=false,
  captionpos=b,
  tabsize=2,
  escapeinside={(*@}{@*)},
  % accents UTF-8 dans le code (listings ne gere pas l'UTF-8 multi-octets)
  literate={é}{{\'e}}1 {è}{{\`e}}1 {ê}{{\^e}}1 {à}{{\`a}}1 {â}{{\^a}}1
           {ù}{{\`u}}1 {û}{{\^u}}1 {ô}{{\^o}}1 {î}{{\^i}}1 {ç}{{\c c}}1
           {É}{{\'E}}1 {À}{{\`A}}1 {€}{{\euro{}}}1 {…}{{\ldots}}1,
}
\lstset{style=r402code}

% ---------------------------------------------------------------
% BOITES : exercice, code, information
% ---------------------------------------------------------------
\tcbset{
  exobox/.style={colback=qapanellight, colframe=qaaccentdark, boxrule=0.6pt,
    arc=2pt, left=8pt, right=8pt, top=6pt, bottom=6pt,
    coltitle=white, colbacktitle=qaaccentdark, fonttitle=\bfseries,
    boxsep=2pt, breakable},
  infobox/.style={colback=white, colframe=qaaccent2dark, boxrule=0.6pt,
    leftrule=3pt, arc=0pt, left=8pt, right=8pt, top=5pt, bottom=5pt,
    coltitle=qaaccent2dark, colbacktitle=white, fonttitle=\bfseries,
    breakable},
  codebox/.style={colback=qatextdark, colframe=qalinelight, boxrule=0.4pt,
    arc=2pt, left=8pt, right=8pt, top=5pt, bottom=5pt, listing only,
    listing options={basicstyle=\ttfamily\small\color{white},
      breaklines=true, showstringspaces=false, columns=fullflexible,
      keepspaces=true, frame=none, numbers=none, backgroundcolor={}}, breakable}
}
\newtcolorbox[auto counter]{exobox}[2][]{exobox,
  title={Exercice~\thetcbcounter\ifstrempty{#2}{}{ --- #2}}, #1}
\newtcolorbox{infobox}[1][\`A retenir]{infobox, title={#1}}
\newtcblisting{codebox}{codebox}

% ---------------------------------------------------------------
% EN-TETE DE PAGE (fil d'Ariane + pagination)
% ---------------------------------------------------------------
\newcommand{\qaheadtitlevalue}{}
\newcommand{\qaheadtitle}[1]{\renewcommand{\qaheadtitlevalue}{#1}}

\pagestyle{fancy}
\fancyhf{}
\renewcommand{\headrulewidth}{0.4pt}
\fancyhead[L]{\small\color{qamuteddark}R4.02 --- \qaheadtitlevalue}
\fancyhead[R]{\small\color{qamuteddark}\thepage\,/\,\pageref{LastPage}}
\fancyfoot[C]{\footnotesize\color{qamuteddark}IUT de Vannes --- Universit\'e Bretagne Sud}
\fancypagestyle{plain}{\fancyhf{}\renewcommand{\headrulewidth}{0pt}%
  \fancyfoot[C]{\footnotesize\color{qamuteddark}IUT de Vannes --- Universit\'e Bretagne Sud}}

% ---------------------------------------------------------------
% BANDEAU DE TITRE (pas une page dediee, pour rester compact)
% \rtpheader{numero ou intitule court}{titre du TP}
% ---------------------------------------------------------------
\newcommand{\ryear}{2026/2027}
\newcommand{\rtpheader}[2]{%
  \thispagestyle{plain}%
  \begin{minipage}[t]{0.68\textwidth}
    {\small\color{qamuteddark}
      Universit\'e Bretagne Sud\\
      IUT de Vannes\\
      D\'epartement Informatique --- \ryear}
  \end{minipage}%
  \hfill
  \begin{minipage}[t]{0.28\textwidth}
    \raggedleft
    \includegraphics[width=2.6cm]{iut_logo.jpg}
  \end{minipage}

  \vspace{6mm}
  \begin{center}
    {\large\bfseries\color{qatextdark} BUT 2 --- R4.02. Qualit\'e de d\'eveloppement}\\[2mm]
    {\normalsize\color{qamuteddark} TP #1}\\[1mm]
    {\Large\bfseries\color{qaaccentdark} #2}\\[2mm]
    \raimention
  \end{center}
  \vspace{4mm}
  \qaheadtitle{#2}
}

%   \begin{document}
%   \rtpheader{1}{Titre du TP}
% ---------------------------------------------------------------
\usepackage[T1]{fontenc}
\usepackage[utf8]{inputenc}
\usepackage[french]{babel}
\usepackage[scale=0.92]{tgheros}
\usepackage{tgpagella}
\usepackage[a4paper, margin=2.2cm, headheight=14pt]{geometry}
\usepackage{graphicx}
\usepackage{xcolor}
\usepackage{listings}
\usepackage[skins,listings,breakable]{tcolorbox}
\usepackage{enumitem}
\usepackage{booktabs}
\usepackage{array}
\usepackage{multirow}
\usepackage{tikz}
\usetikzlibrary{shapes, shapes.geometric, arrows, arrows.meta, positioning, calc}
\usepackage[colorlinks=true, linkcolor=qaaccentdark, urlcolor=qaaccentdark]{hyperref}
\usepackage{fancyhdr}
\usepackage{lastpage}
\usepackage{titlesec}
\usepackage{etoolbox}
\usepackage{eurosym}

% ---------------------------------------------------------------
% MENTION D'ELABORATION (discrete : gris clair, petit corps)
% ---------------------------------------------------------------
\definecolor{qaaimention}{HTML}{A3AAB5}
\newcommand{\raimention}{{\scriptsize\color{qaaimention}Supports \'elabor\'es avec l'aide de Claude Code Opus 5.5}}

% ---------------------------------------------------------------
% PALETTE -- version imprimable des couleurs des diapositives
% ---------------------------------------------------------------
\definecolor{qaaccentdark}{HTML}{7A1680}   % violet du logo IUT, fonce pour rester lisible sur blanc
\definecolor{qaaccent2dark}{HTML}{9B4B94}  % violet clair du logo IUT (accent secondaire)
\definecolor{qapanellight}{HTML}{F4F6F8}   % fond de bloc tres clair
\definecolor{qalinelight}{HTML}{D8DEE4}    % filets
\definecolor{qatextdark}{HTML}{1A1D23}     % texte principal
\definecolor{qamuteddark}{HTML}{5B6472}    % texte secondaire
\definecolor{qacodekwl}{HTML}{7A1680}      % code : mots-cles
\definecolor{qacodestrl}{HTML}{2E7D32}     % code : chaines
\definecolor{qacodecoml}{HTML}{6B7280}     % code : commentaires
\definecolor{qacodeannl}{HTML}{B26A00}     % code : annotations / javadoc

\color{qatextdark}

\setlength{\parskip}{4pt}
\setlength{\parindent}{0pt}

% ---------------------------------------------------------------
% TITRES DE SECTION -- texte blanc sur bande de couleur pleine largeur
% (gere aussi les variantes etoilees \section*)
% ---------------------------------------------------------------
\newcommand{\qasectionband}[2]{%
  \noindent\colorbox{#1}{%
    \parbox[c]{\dimexpr\linewidth-2\fboxsep\relax}{%
      \color{white}#2%
    }%
  }%
}
\newcommand{\secformat}[1]{\qasectionband{qaaccentdark}{\bfseries\Large\thesection\quad #1}}
\newcommand{\secformatnn}[1]{\qasectionband{qaaccentdark}{\bfseries\Large #1}}
\newcommand{\subsecformat}[1]{\qasectionband{qaaccent2dark}{\bfseries\large\thesubsection\quad #1}}
\newcommand{\subsecformatnn}[1]{\qasectionband{qaaccent2dark}{\bfseries\large #1}}

\titleformat{\section}{\normalfont}{}{0pt}{\secformat}
\titleformat{name=\section,numberless}{\normalfont}{}{0pt}{\secformatnn}
\titleformat{\subsection}{\normalfont}{}{0pt}{\subsecformat}
\titleformat{name=\subsection,numberless}{\normalfont}{}{0pt}{\subsecformatnn}
\titleformat{\subsubsection}{\normalfont\bfseries\color{qaaccentdark}}{\thesubsubsection}{0.6em}{}
\titlespacing*{\section}{0pt}{14pt}{8pt}
\titlespacing*{\subsection}{0pt}{10pt}{6pt}
\titlespacing*{\subsubsection}{0pt}{8pt}{4pt}

% ---------------------------------------------------------------
% LISTINGS -- theme clair (lstlisting reste utilisable tel quel)
% \lstcode : police + couleur de base, a reutiliser dans les surcharges
% ---------------------------------------------------------------
\newcommand{\lstcode}{\ttfamily\color{qatextdark}}

\lstdefinelanguage{TypeScript}{
  keywords={typeof, new, true, false, catch, function, return, null, switch, var,
    if, in, while, do, else, case, break, const, let, interface, class, extends,
    implements, enum, export, import, this, super, constructor, public, private,
    protected, static, async, await, throw, throws, try, finally, void},
  ndkeywords={boolean, any, number, string, int, long, double},
  sensitive=true,
  comment=[l]{//},
  morecomment=[s]{/*}{*/},
  morestring=[b]',
  morestring=[b]"
}

\lstdefinestyle{r402code}{
  basicstyle=\lstcode\footnotesize,
  keywordstyle=\color{qacodekwl}\bfseries,
  ndkeywordstyle=\color{qaaccent2dark}\bfseries,
  stringstyle=\color{qacodestrl},
  commentstyle=\color{qacodecoml}\itshape,
  morecomment=[s][\color{qacodeannl}]{/**}{*/},
  backgroundcolor=\color{qapanellight},
  rulecolor=\color{qalinelight},
  frame=single,
  framesep=4pt,
  xleftmargin=6pt,
  xrightmargin=6pt,
  numbers=left,
  numberstyle=\tiny\color{qamuteddark},
  numbersep=8pt,
  keepspaces=true,
  columns=fullflexible,
  showstringspaces=false,
  showspaces=false,
  showtabs=false,
  breaklines=true,
  breakatwhitespace=false,
  captionpos=b,
  tabsize=2,
  escapeinside={(*@}{@*)},
  % accents UTF-8 dans le code (listings ne gere pas l'UTF-8 multi-octets)
  literate={é}{{\'e}}1 {è}{{\`e}}1 {ê}{{\^e}}1 {à}{{\`a}}1 {â}{{\^a}}1
           {ù}{{\`u}}1 {û}{{\^u}}1 {ô}{{\^o}}1 {î}{{\^i}}1 {ç}{{\c c}}1
           {É}{{\'E}}1 {À}{{\`A}}1 {€}{{\euro{}}}1 {…}{{\ldots}}1,
}
\lstset{style=r402code}

% ---------------------------------------------------------------
% BOITES : exercice, code, information
% ---------------------------------------------------------------
\tcbset{
  exobox/.style={colback=qapanellight, colframe=qaaccentdark, boxrule=0.6pt,
    arc=2pt, left=8pt, right=8pt, top=6pt, bottom=6pt,
    coltitle=white, colbacktitle=qaaccentdark, fonttitle=\bfseries,
    boxsep=2pt, breakable},
  infobox/.style={colback=white, colframe=qaaccent2dark, boxrule=0.6pt,
    leftrule=3pt, arc=0pt, left=8pt, right=8pt, top=5pt, bottom=5pt,
    coltitle=qaaccent2dark, colbacktitle=white, fonttitle=\bfseries,
    breakable},
  codebox/.style={colback=qatextdark, colframe=qalinelight, boxrule=0.4pt,
    arc=2pt, left=8pt, right=8pt, top=5pt, bottom=5pt, listing only,
    listing options={basicstyle=\ttfamily\small\color{white},
      breaklines=true, showstringspaces=false, columns=fullflexible,
      keepspaces=true, frame=none, numbers=none, backgroundcolor={}}, breakable}
}
\newtcolorbox[auto counter]{exobox}[2][]{exobox,
  title={Exercice~\thetcbcounter\ifstrempty{#2}{}{ --- #2}}, #1}
\newtcolorbox{infobox}[1][\`A retenir]{infobox, title={#1}}
\newtcblisting{codebox}{codebox}

% ---------------------------------------------------------------
% EN-TETE DE PAGE (fil d'Ariane + pagination)
% ---------------------------------------------------------------
\newcommand{\qaheadtitlevalue}{}
\newcommand{\qaheadtitle}[1]{\renewcommand{\qaheadtitlevalue}{#1}}

\pagestyle{fancy}
\fancyhf{}
\renewcommand{\headrulewidth}{0.4pt}
\fancyhead[L]{\small\color{qamuteddark}R4.02 --- \qaheadtitlevalue}
\fancyhead[R]{\small\color{qamuteddark}\thepage\,/\,\pageref{LastPage}}
\fancyfoot[C]{\footnotesize\color{qamuteddark}IUT de Vannes --- Universit\'e Bretagne Sud}
\fancypagestyle{plain}{\fancyhf{}\renewcommand{\headrulewidth}{0pt}%
  \fancyfoot[C]{\footnotesize\color{qamuteddark}IUT de Vannes --- Universit\'e Bretagne Sud}}

% ---------------------------------------------------------------
% BANDEAU DE TITRE (pas une page dediee, pour rester compact)
% \rtpheader{numero ou intitule court}{titre du TP}
% ---------------------------------------------------------------
\newcommand{\ryear}{2026/2027}
\newcommand{\rtpheader}[2]{%
  \thispagestyle{plain}%
  \begin{minipage}[t]{0.68\textwidth}
    {\small\color{qamuteddark}
      Universit\'e Bretagne Sud\\
      IUT de Vannes\\
      D\'epartement Informatique --- \ryear}
  \end{minipage}%
  \hfill
  \begin{minipage}[t]{0.28\textwidth}
    \raggedleft
    \includegraphics[width=2.6cm]{iut_logo.jpg}
  \end{minipage}

  \vspace{6mm}
  \begin{center}
    {\large\bfseries\color{qatextdark} BUT 2 --- R4.02. Qualit\'e de d\'eveloppement}\\[2mm]
    {\normalsize\color{qamuteddark} TP #1}\\[1mm]
    {\Large\bfseries\color{qaaccentdark} #2}\\[2mm]
    \raimention
  \end{center}
  \vspace{4mm}
  \qaheadtitle{#2}
}

\graphicspath{{../../latex/common/assets/}{./}}

% A ajuster avant diffusion
\newcommand{\rdatelimite}{la veille de la séance du TP6, à 23\,h\,59}

\begin{document}

\rtpheader{5}{Tests de bout en bout de l'interface web avec Playwright}

\begin{infobox}[Objectifs]
\begin{itemize}[leftmargin=1.2em, itemsep=1pt]
  \item Piloter un vrai navigateur depuis des tests JUnit avec Playwright for Java.
  \item Partir d'un test généré par \texttt{codegen} et le transformer en test robuste (sélecteurs, attentes, Page Object).
  \item Trouver des bugs qu'aucun test côté serveur ne pouvait voir.
  \item Repérer et corriger un test instable ; déboguer avec le \textit{trace viewer}.
  \item Exécuter un scénario Cucumber du TP2 à travers l'interface.
\end{itemize}
\end{infobox}

Continuez dans le dépôt de votre binôme, à partir de votre rendu du TP4. Docker doit être démarré : les tests de bout en bout utilisent la même base PostgreSQL de test que le TP4.

% ================================================================
\section{Mise en place (10 min)}

Le projet de départ contient :
\begin{itemize}[leftmargin=1.2em, itemsep=1pt]
  \item le front : une page de connexion (\texttt{connexion.html}) et une application Vue 3 (\texttt{index.html} et \texttt{js/}), dans \texttt{src/main/resources/static}. Vue est utilisé \textbf{sans étape de construction} : pas de Node.js, une modification d'un composant (\texttt{js/composants/*.js}) est prise en compte dès que l'application est relancée ;
  \item la dépendance Playwright for Java, et le plugin Maven qui donne accès à ses outils (\texttt{./mvnw exec:java -Dexec.args="..."}) ;
  \item la classe \texttt{BaseE2ETest} (paquetage de test \texttt{fr.iutvannes.banque.e2e}), qui démarre l'application sur un port aléatoire avec la base de test, puis un navigateur Chromium, et enregistre une trace par test ;
  \item deux utilisateurs de test : \texttt{alice} / \texttt{alice123} (comptes \texttt{FR76-001} et \texttt{FR76-002}) et \texttt{bruno} / \texttt{bruno123} (compte \texttt{FR76-003}).
\end{itemize}

\begin{exobox}{Prise en main}
\begin{enumerate}[leftmargin=1.4em]
  \item Téléchargez le navigateur utilisé par Playwright (une seule fois) :
\begin{codebox}
./mvnw exec:java -Dexec.args="install chromium"
\end{codebox}
  \item Lancez l'application (\texttt{docker compose up -d} puis \texttt{./mvnw spring-boot:run}) et utilisez-la à la main : connexion, consultation des comptes, virement. Notez dans \texttt{docs/tp5.md} les deux ou trois parcours qui vous semblent \textbf{critiques} pour la banque.
  \item Lisez \texttt{BaseE2ETest}. Dans \texttt{docs/tp5.md}, expliquez en une phrase chacun : pourquoi un \texttt{BrowserContext} neuf par test ? Où et quand les données de la base sont-elles remises à zéro ? À quoi sert l'option \texttt{-Dheaded=true} ?
\end{enumerate}
\end{exobox}

% ================================================================
\section{Partir d'un test enregistré (15 min)}

\begin{exobox}{Le générateur de code}
\begin{enumerate}[leftmargin=1.4em]
  \item L'application étant lancée, démarrez le générateur :
\begin{codebox}
./mvnw exec:java -Dexec.args="codegen http://localhost:8080"
\end{codebox}
  \item Dans le navigateur qui s'ouvre, connectez-vous avec \texttt{alice}, faites un virement de 150\,\euro{} de \texttt{FR76-001} vers \texttt{FR76-002}, puis revenez à la liste des comptes. Dans la fenêtre de l'inspecteur, choisissez la cible « Java / JUnit ».
  \item Copiez le code généré dans une classe \texttt{VirementEnregistreTest} qui hérite de \texttt{BaseE2ETest}, et adaptez-le pour qu'il tourne (URL, objet \texttt{page}). Lancez-le, en mode visible puis invisible.
  \item Critiquez ce code dans \texttt{docs/tp5.md} : quels sélecteurs sont fragiles ? Que vérifie-t-il réellement ? Que se passerait-il si on renommait le bouton « Valider » en « Confirmer » ? Si on ajoutait un champ au formulaire ?
\end{enumerate}
\end{exobox}

% ================================================================
\section{Des tests robustes (30 min)}

\begin{exobox}{Page Objects}
Dans le paquetage \texttt{fr.iutvannes.banque.e2e.pages}, écrivez un Page Object par écran : \texttt{PageConnexion}, \texttt{PageComptes}, \texttt{PageVirement}. Chacun expose des \textbf{actions métier} (\texttt{seConnecter(identifiant, motDePasse)}, \texttt{virer(source, cible, montant)}, \texttt{solde(numero)}\ldots) et ne contient aucune assertion.
\begin{itemize}[leftmargin=1.2em]
  \item Utilisez en priorité \texttt{getByRole} et \texttt{getByLabel}, puis \texttt{getByText}.
  \item Quand aucun de ces sélecteurs ne convient (par exemple la cellule du solde d'un compte dans un tableau), ajoutez un attribut \texttt{data-testid} dans le composant Vue concerné (\texttt{src/main/resources/static/js/composants}), par exemple \texttt{:data-testid="'solde-' + compte.numero"}.
  \item Si \texttt{getByRole} ou \texttt{getByLabel} ne trouve pas un élément pourtant visible, regardez le HTML : un champ sans \texttt{<label>}, un \texttt{<div>} cliquable au lieu d'un \texttt{<button>} ? Corrigez le composant Vue plutôt que de contourner avec un sélecteur CSS. Notez ces corrections dans \texttt{docs/tp5.md} : ce sont aussi des corrections d'\textbf{accessibilité}.
\end{itemize}
\end{exobox}

\begin{exobox}{Les parcours critiques}
Réécrivez le test enregistré, puis ajoutez les autres parcours, dans une classe \texttt{VirementE2ETest}. Toutes les vérifications passent par \texttt{assertThat(locator)} de Playwright ; aucun \texttt{Thread.sleep}, aucun \texttt{textContent()} suivi d'un \texttt{assertEquals}.
\begin{enumerate}[leftmargin=1.4em]
  \item Connexion réussie (le nom du client est affiché), connexion refusée (message d'erreur, on reste sur la page de connexion).
  \item Virement accepté : message de confirmation \textbf{et} nouveaux soldes affichés dans la liste des comptes.
  \item Virement de \textbf{150,50\,\euro{}} : les soldes affichés sont-ils corrects au centime près ?
  \item Virement refusé pour solde insuffisant (connecté en \texttt{bruno}) : le message affiché à l'utilisateur est-il compréhensible ?
  \item Virement de plus de 1\,000\,\euro{} : une fenêtre de confirmation (\texttt{role="dialog"}) s'affiche ; testez la confirmation \textbf{et} l'annulation (aucun solde ne change).
\end{enumerate}
Au moins deux de ces parcours révèlent un bug \textbf{du front}. Pour chacun : notez-le dans \texttt{docs/tp5.md}, expliquez pourquoi les tests des TP1 à TP4 ne pouvaient pas le voir, corrigez le composant Vue, relancez.
\end{exobox}

% ================================================================
\section{Instabilité et débogage (10 min)}

\begin{exobox}{Un test instable ?}
\begin{enumerate}[leftmargin=1.4em]
  \item Remplacez temporairement \texttt{@Test} par \texttt{@RepeatedTest(10)} sur deux de vos tests et lancez-les. Passent-ils dix fois sur dix ?
  \item Écrivez volontairement une version fragile d'un test : lire le message avec \texttt{textContent()} juste après le clic, et le comparer avec \texttt{assertEquals}. Répétez-la dix fois. Expliquez le résultat.
  \item Ouvrez la trace d'un échec :
\begin{codebox}
./mvnw exec:java -Dexec.args="show-trace target/traces/<nom-du-test>.zip"
\end{codebox}
  Parcourez les actions une à une. À quel moment la page affichait-elle le message ? Qu'y avait-il dans l'onglet \textit{Network} au moment de la lecture ? Supprimez ensuite la version fragile.
\end{enumerate}
\end{exobox}

% ================================================================
\section{Le scénario Cucumber, dans le navigateur (15 min)}

\begin{exobox}{Mêmes scénarios, autres étapes}
Au TP2, vos étapes Cucumber appelaient directement \texttt{VirementService}. On veut maintenant exécuter \textbf{un} scénario de \texttt{virement.feature} à travers l'interface, sans modifier le texte du scénario.
\begin{enumerate}[leftmargin=1.4em]
  \item Marquez ce scénario avec un tag \texttt{@ui}.
  \item Créez un second lanceur \texttt{CucumberUiTest}, avec un paquetage d'étapes distinct, \textbf{qui ne soit pas un sous-paquetage} de \texttt{acceptation} (par exemple \texttt{fr.iutvannes.banque.e2e.cucumber}), et un filtre sur le tag :\\
  \texttt{@ConfigurationParameter(key = FILTER\_TAGS\_PROPERTY\_NAME, value = "@ui")}.
  \item Écrivez les étapes de ce paquetage en utilisant vos Page Objects. Les étapes « Étant donné un compte\ldots » doivent préparer les données \textbf{sans passer par l'interface} (directement en base, comme dans \texttt{BaseE2ETest}).
  \item Le lanceur du TP2 exécute-t-il toujours tous les scénarios, y compris \texttt{@ui} ? Est-ce souhaitable ? Justifiez en deux phrases dans \texttt{docs/tp5.md}.
\end{enumerate}
\textit{Indice} : pour que les étapes aient accès à l'application démarrée et à son port, annotez une classe du paquetage d'étapes avec \texttt{@CucumberContextConfiguration} et \texttt{@SpringBootTest(webEnvironment = RANDOM\_PORT)}, comme \texttt{ConfigurationCucumber} le fait pour le TP2. Le port est disponible par \texttt{@Value("\$\{local.server.port\}")}. Le navigateur se lance dans des méthodes \texttt{@BeforeAll} / \texttt{@AfterAll} statiques de Cucumber (\texttt{io.cucumber.java}), un onglet par scénario dans \texttt{@Before} / \texttt{@After}.
\end{exobox}

% ================================================================
\section{Rendu}

\begin{infobox}[À rendre avant \rdatelimite]
Sur le dépôt du binôme, un \textbf{tag} \texttt{tp5} pointant sur le commit rendu :
\begin{codebox}
git tag tp5
git push origin tp5
\end{codebox}
Le commit doit contenir :
\begin{itemize}[leftmargin=1.2em, itemsep=1pt]
  \item les Page Objects et \texttt{VirementE2ETest} (le test enregistré \texttt{VirementEnregistreTest} peut être supprimé) ;
  \item les composants Vue corrigés (bugs, \texttt{data-testid}, accessibilité) ;
  \item \texttt{CucumberUiTest} et ses étapes ;
  \item \texttt{docs/tp5.md} : parcours critiques, réponses sur \texttt{BaseE2ETest}, critique du code généré, journal des bugs du front, analyse du test instable, choix pour les scénarios \texttt{@ui} ;
  \item une suite complète qui passe : \texttt{./mvnw test} (Docker démarré).
\end{itemize}
\end{infobox}

\end{document}
